\begin{abstract}
    Transport models routinely anchor points of interest (POIs) at
    the centroid of their building or parcel. On large multi-building
    sites---university campuses, hospitals, malls, industrial
    parks---this can misrepresent pedestrian access by hundreds of
    metres, distorting first-/last-mile transit access, wheelchair
    routing, and curbside drop-off precision for emerging autonomous
    fleets. We quantify how explicit \texttt{entrance=*} detail in
    OpenStreetMap (OSM) shifts accessibility-style metrics relative
    to centroid proxies. The study
    combines a global, population/built-volume-weighted random sample
    with a Qu\'ebec validation against street-level imagery,
    authoritative registers, and field-surveyed entrances. We release
    the full toolchain---a Rust/React sampling and review service
    over Overpass and PostGIS, with a MapLibre interface that
    persists typed geodesic measurements between POI / entrance /
    centroid anchors and the nearest transit, walking, cycling,
    parking, and driving targets---together with the tabulated
    outputs.
\end{abstract}